\documentclass[11pt]{article}

\usepackage[T1]{fontenc}
\usepackage{amsmath,amssymb,amsthm}
\usepackage{booktabs,tabularx}
\usepackage[hidelinks]{hyperref}
\usepackage[margin=1in]{geometry}

\newtheorem{theorem}{Theorem}[section]
\newtheorem{lemma}[theorem]{Lemma}
\newtheorem{proposition}[theorem]{Proposition}
\newtheorem{corollary}[theorem]{Corollary}
\theoremstyle{definition}
\newtheorem{definition}[theorem]{Definition}
\theoremstyle{remark}
\newtheorem{remark}[theorem]{Remark}

\newcommand{\En}{\mathcal E}
\newcommand{\Sn}{S}
\newcommand{\dist}{d_t}
\newcommand{\Del}{\operatorname{del}}
\newcommand{\Child}{\operatorname{Child}}
\newcommand{\Fix}{\operatorname{Fix}}
\newcommand{\id}{\iota}
\newcommand{\cycle}[1]{(#1)}

\title{Propagation of Extremal Classes in Sorting by Transpositions:\
  a Route to the Diameter of $S_{18}$}
\author{Luiz A. G. Silva \and Luis A. B. Kowada}
\date{}

\begin{document}
\maketitle

\begin{abstract}
Sorting by transpositions asks for the minimum number of exchanges of two
adjacent blocks that transform a permutation into the identity.  We develop a
class-level reduction for determining the transposition diameter at order
$18$.  Deleting any symbol changes transposition distance by at most one.
Consequently, a permutation whose distance is one larger than the diameter at
the preceding order must have every one-symbol deletion in the preceding
extremal set.  We formulate this reduction for extended-toric classes of
cyclic-target pairs, prove that one representative has only quadratically many
child insertions before canonicalization, and obtain a finite pipeline
$\En_{15}\to\En_{16}\to\En_{17}\to S_{18}$.  The first two stages compute the
complete extremal sets at orders $15$ and $16$; the next stage propagates and
classifies the surviving order-$17$ classes; the last stage requires a sorting
sequence of length at most ten for every surviving order-$18$ class.  Together
with the known lower bound, successful exhaustive generation, search, and
verification would establish $TD(18)=10$.  This article records the reduction,
the proof obligations, and the certificate design.  The large computation is
in progress, and no completed $TD(18)$ result is claimed in this version.
\end{abstract}

\paragraph{Keywords.}
sorting by transpositions; transposition diameter; toric equivalence;\newline
computer-assisted proof; exhaustive generation

\paragraph{MSC 2020.}
05A05, 05A15, 20B30, 68R05.

\section{Introduction}\label{sec:introduction}

Let $[n]:=\{1,\ldots,n\}$ and
$\pi=[\pi_1\ \cdots\ \pi_n]\in S_n$.  A \emph{transposition}
exchanges two adjacent, nonempty blocks while preserving the order inside
each block.  The \emph{transposition distance} $\dist(\pi)$ is the minimum
number of transpositions that sort $\pi$, and the \emph{transposition
diameter} is
\[
  TD(n):=\max_{\pi\in S_n}\dist(\pi).
\]
Bafna and Pevzner introduced this problem and its cycle-graph lower
bound~\cite{BafnaPevzner1998}.  Eriksson et al. determined the exact diameter
through order $15$ and proved the recurrence
$TD(n)\leq TD(n-3)+2$ for $n\geq9$~\cite{ErikssonEtAl2001}.
The twisted-bracelet computation of Silva et al. subsequently established
$TD(16)=9$~\cite{SilvaEtAl2026}.  The known values relevant here are therefore
\[
  TD(14)=8,\qquad TD(15)=9,\qquad TD(16)=9,
  \qquad TD(17)=10.
\]
The last equality also follows because the lower bound of Elias and
Hartman~\cite{EliasHartman2006} and the upper bound of Eriksson et al. both
equal ten at order $17$; see also Korchmaros~\cite{Korchmaros2015}.

At order $18$, the reverse-permutation lower bound gives $TD(18)\geq10$
\cite{MeidanisWalterDias1997}, whereas the recurrence and $TD(15)=9$ give
$TD(18)\leq11$.  Thus it remains to rule out distance eleven.  Direct
enumeration of $S_{18}$ is not a credible approach.  Even after quotienting
by the usual distance-preserving symmetries, the complete fixed-point-free
space at the cyclic order needed for $S_{16}$ is expected to contain about
$2.1\cdot10^{11}$ classes.  The central purpose of this paper is to replace
enumeration at orders $17$ and $18$ by propagation from the much smaller
sets of classes that are already extremal at the preceding order.

The reduction rests on an elementary observation.  If a symbol is deleted
from a permutation, any sorting sequence for the shorter permutation can be
lifted while carrying the deleted symbol with an adjacent block.  At most one
additional transposition then puts that symbol in its final position.  Hence a
new extremal permutation can occur only when all of its deletions were already
extremal.  At class level, deletion is compatible with toric rotation and
reflection, and all child classes of a representative are obtained from only
$N^2$ insertions on its $N$-symbol cyclic domain.

Our contributions in this working draft are as follows.
\begin{enumerate}
  \item We prove the deletion inequality and its extremal-set consequence.
  \item We formulate deletion and insertion for cyclic-target pairs and state
        the compatibility needed to work directly with extended-toric classes.
  \item We derive a four-stage exhaustive computation from $\En_{15}$ to the
        order-$18$ candidates and isolate the exact proof obligation at each
        stage.
  \item We specify persistent class sets, sorting witnesses, failure checks,
        and deterministic manifests sufficient to audit the computation.
\end{enumerate}
The present implementation realizes these stages, but the large run and its
checked artifacts are not part of this draft.  Section~\ref{sec:preliminaries}
sets the notation.  Section~\ref{sec:deletion} proves the deletion reduction.
Section~\ref{sec:classes} transfers it to extended-toric classes.
Section~\ref{sec:pipeline} gives the four stages and the conditional diameter
theorem.  Section~\ref{sec:certificate} describes the certificate and
verification obligations, and Section~\ref{sec:status} records the remaining
work.

\section{Preliminaries}\label{sec:preliminaries}

For $1\leq i<j<k\leq n+1$, let $\tau(i,j,k)$ denote the transposition
that sends
\[
 [\pi_1\cdots\pi_{i-1}\mid
   \pi_i\cdots\pi_{j-1}\mid
   \pi_j\cdots\pi_{k-1}\mid
   \pi_k\cdots\pi_n]
\]
to
\[
 [\pi_1\cdots\pi_{i-1}\mid
   \pi_j\cdots\pi_{k-1}\mid
   \pi_i\cdots\pi_{j-1}\mid
   \pi_k\cdots\pi_n].
\]
For $x\in[n]$, define the order-preserving standardization map
$s_x:[n]\setminus\{x\}\to[n-1]$ by
\[
  s_x(y):=
  \begin{cases}
    y, & y<x,\\
    y-1, & y>x.
  \end{cases}
\]
If $x=\pi_r$, the \emph{deletion of $x$ from $\pi$} is
\begin{equation}\label{eq:symbol-deletion}
  \Del_x(\pi):=
  [s_x(\pi_1)\ \cdots\ s_x(\pi_{r-1})\
   \ s_x(\pi_{r+1})\ \cdots\ s_x(\pi_n)]\in S_{n-1}.
\end{equation}
Thus deletion removes the entry $x$ and standardizes the remaining alphabet.
This order-preserving relabeling does not affect transposition distance.

We use the cyclic-target model developed for the twisted-bracelet computation
of $TD(16)$~\cite{SilvaEtAl2026}.  Let $E_N=\{0,\ldots,N-1\}$.  A
\emph{cyclic-target pair} is a pair $(\omega,\beta)$ such that $\beta$ and
$\rho:=\omega\beta$ are $N$-cycles.  We call $\beta$ the current cycle,
$\rho$ the target cycle, and $\omega=\rho\beta^{-1}$ the algebraic
permutation.  An ordinary permutation in $S_n$ is represented on the cyclic
domain of size $N=n+1$.  Renaming the target cycle to
$\cycle{0\ 1\ \cdots\ N-1}$ recovers the ordinary sorting problem and
preserves distance.

Toric rotations and reflection generate a dihedral action by
distance-preserving maps.  Its orbits are the \emph{extended-toric classes}.
For fixed-point-free algebraic permutations of prescribed annotated cycle
type, realizable twisted bracelets correspond bijectively to these
classes~\cite{SilvaEtAl2026}.  We use that correspondence only for exhaustive
generation; the deletion reduction below applies to classes with or without
fixed points.

\begin{definition}[Extremal classes]
For an order $n$ whose diameter is known, let $\En_n$ be the set of
extended-toric classes represented by permutations $\pi\in S_n$ satisfying
$\dist(\pi)=TD(n)$.
\end{definition}

Distance is constant on an extended-toric class, so this definition does not
depend on the representative.  The distinction between permutations and
classes remains important: a class table stores one canonical representative,
whereas symbol deletion is evaluated through all circular positions of a
cyclic-target representative.

\section{The deletion reduction}\label{sec:deletion}

\begin{lemma}[One-symbol deletion]\label{lem:deletion}
For every permutation $\pi$ and every symbol $x$ of $\pi$,
\[
  \dist(\pi)\leq \dist(\Del_x(\pi))+1.
\]
\end{lemma}

\begin{proof}
Let $\tau_1,\ldots,\tau_s$ be a shortest sorting sequence for
$\Del_x(\pi)$.
At every step, lift the two exchanged intervals to the current word containing
$x$.  If $x$ lies in a gap touched by the exchange, assign it to either
neighboring interval that contains that gap.  Deleting $x$ after the lifted
move therefore gives exactly the move on the shorter word.  Inductively, the
lifted sequence transforms $\pi$ into a word whose deletion of $x$ is the
identity on the other $n-1$ symbols.

The resulting word is the identity with $x$ inserted in one of its gaps.  If
$x$ is not already in its identity position, exchange the singleton block
containing $x$ with the adjacent interval between its current and prescribed
positions.  This is one transposition and completes the sorting.  Thus
$\dist(\pi)\leq s+1$.
\end{proof}

\begin{corollary}[All removals are extremal]\label{cor:all-removals}
Suppose that $TD(n)=TD(n-1)+1$.  If
$\pi\in S_n$ satisfies $\dist(\pi)=TD(n)$, then
$\Del_x(\pi)\in\En_{n-1}$ for every symbol $x$, where membership is understood
up to extended-toric equivalence.
\end{corollary}

\begin{proof}
Lemma~\ref{lem:deletion} gives
\[
 TD(n)=\dist(\pi)\leq\dist(\Del_x(\pi))+1
 \leq TD(n-1)+1=TD(n).
\]
Both inequalities are equalities, and hence
$\dist(\Del_x(\pi))=TD(n-1)$.
\end{proof}

The contrapositive is the useful computational form.  If one deletion of an
order-$n$ instance lies outside $\En_{n-1}$, then the instance cannot be one
move beyond $TD(n-1)$.  This is a membership test, not a search: after the
preceding extremal set has been materialized, every deletion is resolved by a
canonical key lookup.

\section{Deletion and insertion on classes}\label{sec:classes}

Let $(\omega,\beta)$ be a cyclic-target pair on $E_N$, with
$\rho=\omega\beta$.  For $x\in E_N$, delete $x$ from the cyclic notation of
both full cycles $\beta$ and $\rho$, close the two cycles, and apply the
unique order-preserving bijection from $E_N\setminus\{x\}$ to $E_{N-1}$.
This is the zero-based analogue of the standardization in
\eqref{eq:symbol-deletion}.  We denote the resulting full cycles by
$\Del_x(\beta)$ and $\Del_x(\rho)$.  They define the smaller pair
\[
  \Del_x(\omega,\beta)
  :=\bigl(\Del_x(\rho)\Del_x(\beta)^{-1},\,\Del_x(\beta)\bigr).
\]
This simultaneous deletion is essential.  Deleting a symbol only from a
stored one-line permutation does not preserve the cyclic target.

The inverse construction inserts the new largest symbol $N$ into one gap of
$\beta$ and independently into one gap of $\rho$.  There are $N$ gaps in
each cycle and therefore $N^2$ raw insertions.  The algebraic permutation of
the resulting pair is reconstructed as $\rho'({\beta'})^{-1}$.

\begin{proposition}[Class compatibility]\label{prop:class-compatibility}
The multiset
\[
  \bigl\{[\Del_x(\omega,\beta)]:x\in E_N\bigr\}
\]
of extended-toric classes depends only on the extended-toric class of
$(\omega,\beta)$.  Reflection commutes with deletion up to the induced
order-preserving relabeling.  Moreover, the $N^2$ raw insertions of any one
representative, followed by canonicalization on $E_{N+1}$, contain every
extended-toric child class and no class whose deletion cannot return to the
parent class.
\end{proposition}

\begin{proof}
A toric rotation changes only the distinguished origin of each full cycle.
Deleting all possible symbols after this change permutes the resulting
deletions and conjugates each by the induced relabeling.  Reversal changes the
orientation of both full cycles; deleting a position before or after reversal
again differs only by the induced relabeling.  These maps are precisely the
rotation and reflection maps defining extended-toric equivalence, proving the
multiset statement.

For the converse direction, let a child class have a deletion in the parent
class.  Choose the equivalence taking that deletion to the fixed parent
representative and apply the same renaming to the child.  The deleted symbol
then occupies one gap of the representative's current cycle and one gap of
its target cycle.  The renamed child is consequently one of the $N^2$ raw
insertions.  Canonicalization removes repetitions and changes no class.
\end{proof}

Write $\Child(C)$ for the canonical child classes obtained from a class $C$,
and define the complete child set of a class family $A$ by
\[
  \Child(A):=\bigcup_{C\in A}\Child(C).
\]
For a child class $C$, write $\Del(C)$ for its multiset of deletion classes.
Proposition~\ref{prop:class-compatibility} makes both definitions independent
of all representative choices.

An algebraic fixed point of $\omega$ is an adjacency shared by $\beta$ and
$\rho$.  Contracting the corresponding edge deletes the fixed symbol from
both cycles and preserves transposition distance~\cite{SilvaEtAl2026}.
Conversely, inserting a symbol into the same corresponding gap of both cycles
restores a fixed point and preserves distance.  We denote the classes obtained
in this way by $\operatorname{Expand}(A)$.

\begin{proposition}[Fixed-point part of $\En_{16}$]
The classes in $\En_{16}$ with an algebraic fixed point are exactly
$\operatorname{Expand}(\En_{15})$.  Each such class has exactly one algebraic
fixed point.
\end{proposition}

\begin{proof}
Contraction preserves distance.  Hence contracting a fixed point from a class
of distance $TD(16)=9$ produces an order-$15$ class of distance nine, which
belongs to $\En_{15}$.  Conversely, fixed-point expansion of any class in
$\En_{15}$ preserves its distance nine and therefore gives a class in
$\En_{16}$.  Since $TD(15)=9>TD(14)=8$, a member of $\En_{15}$ cannot itself
have an algebraic fixed point.  The restored class consequently has exactly
one.
\end{proof}

\section{The extremal-class pipeline}\label{sec:pipeline}

The cyclic domain for an ordinary order-$n$ permutation has size $N=n+1$.
Thus the four stages below use domains $16,17,18,$ and $19$, respectively.
The phrase \emph{failure through depth $k$} always means exhaustive failure
of an exact bounded search, not failure of a heuristic.  A successful search
is retained as an explicit sequence and proves only an upper bound.

\begin{table}[ht]
\centering
\begin{tabularx}{\textwidth}{@{}ll
  >{\raggedright\arraybackslash}X
  >{\raggedright\arraybackslash}X@{}}
\toprule
Stage & Cyclic domain & Input or enumeration & Required conclusion \\
\midrule
$\En_{15}$ & $16$ & all fixed-point-free classes
  & fail at $8$, sort at $9$ \\
$\En_{16}$ & $17$ & $\operatorname{Expand}(\En_{15})$ and all
  fixed-point-free classes & fail at $8$, sort at $9$ \\
$\En_{17}$ & $18$ & children of $\En_{16}$ with all removals in $\En_{16}$
  & fail at $9$, sort at $10$ \\
$S_{18}$ candidates & $19$ & children of $\En_{17}$ with all removals in
  $\En_{17}$ & sort at $10$ \\
\bottomrule
\end{tabularx}
\caption{The four exhaustive stages.  The last stage need not compute the
exact set $\En_{18}$; it only certifies a ten-move upper bound for every class
that could have larger distance.}
\label{tab:pipeline}
\end{table}

\subsection{Stage 1: computing \texorpdfstring{$\En_{15}$}{E15}}

Since $TD(15)=9>TD(14)=8$, contraction shows that every member of
$\En_{15}$ is fixed-point-free.  We enumerate each realizable
fixed-point-free twisted bracelet on the $16$-symbol cyclic domain, search it
exhaustively through depth eight, and retain precisely the failures.  Every
retained class must also have a checked nine-move sorting sequence.  The two
tests establish equality rather than merely an upper bound.

The stage will be compared with the known hard instances from the $TD(16)$
computation, including its residual depth-eight obstructions and the published
diametral examples.  These comparisons are regression tests; completeness
comes from exhaustive twisted-bracelet generation across every admissible
annotated cycle type.

\subsection{Stage 2: computing \texorpdfstring{$\En_{16}$}{E16}}

The fixed-point part is $\operatorname{Expand}(\En_{15})$ by the proposition
above.  For the remaining part, we enumerate all fixed-point-free classes on
the $17$-symbol cyclic domain and retain exactly those that fail exhaustive
search through depth eight.  The theorem $TD(16)=9$ supplies the ambient upper
bound, and stored nine-move sequences provide local witnesses.

Let $\lVert\pi\rVert_3$ denote the Bafna--Pevzner 3-norm.  The approximation
guarantees used for pruning give the integer upper bound
\[
 U(k):=
 \begin{cases}
   \lfloor 3k/2\rfloor, & 0\leq k<8,\\
   \lfloor 11k/8\rfloor, & k\geq8,
 \end{cases}
 \qquad k=\lVert\pi\rVert_3.
\]
The first line follows from the Bafna--Pevzner $1.5$ guarantee
\cite{BafnaPevzner1998}; the stronger $1.375$ guarantee is applicable only
when $k\geq8$~\cite{EliasHartman2006}.  Consequently, in the depth-eight
classification, 3-norm at most five is discharged by the approximation bound,
whereas 3-norm six must be searched with two units of slack.  Norms seven and
eight have, respectively, one unit and zero units of slack; the latter case
reduces to the existence of an all-2-move sorting sequence.

\subsection{Stage 3: computing \texorpdfstring{$\En_{17}$}{E17}}

Generate $\Child(\En_{16})$, canonicalize on the $18$-symbol domain, and
deduplicate globally.  A child $C$ reaches exact search only if every member of
$\Del(C)$ lies in $\En_{16}$.  Search each survivor through depth nine.
Because $TD(17)=10$, the failures are exactly $\En_{17}$; for each failure we
also store a sorting sequence of length at most ten.

\subsection{Stage 4: certifying order 18}

Generate $\Child(\En_{17})$ on the $19$-symbol cyclic domain and retain a
child $C$ only if every member of $\Del(C)$ lies in $\En_{17}$.  This stage
does not need to determine which survivors have distance exactly ten.  It
only needs a checked sorting sequence of length at most ten for every survivor.

\begin{theorem}[Conditional certification theorem]\label{thm:conditional}
Assume that the first three stages of Table~\ref{tab:pipeline} produce the
complete sets $\En_{15}$, $\En_{16}$, and $\En_{17}$, with the stated
exhaustive failures and sorting witnesses.  If every class surviving the
fourth-stage removal filter has a sorting sequence of length at most ten, then
$TD(18)=10$.
\end{theorem}

\begin{proof}
The reverse-permutation lower bound gives $TD(18)\geq10$.  Suppose for a
contradiction that some $\pi\in S_{18}$ has distance at least eleven.  Since
$TD(17)=10$, Lemma~\ref{lem:deletion} implies
\[
  11\leq\dist(\pi)\leq\dist(\Del_x(\pi))+1\leq11
\]
for every symbol $x$.  Thus every deletion has distance ten and belongs to
$\En_{17}$.  Proposition~\ref{prop:class-compatibility} places the class of
$\pi$ among the children generated in Stage~4, and it survives the
all-removals filter.  By hypothesis it sorts within ten moves, contradicting
$\dist(\pi)\geq11$.  Hence $TD(18)\leq10$, and equality follows.
\end{proof}

\begin{remark}
The theorem separates mathematical completeness from search success.  A
ten-move witness for every class in an incomplete candidate list proves
nothing about omitted classes.  Conversely, exhaustive candidate generation
without replayable witnesses does not prove the upper bound.
\end{remark}

\section{Certificate and verification}\label{sec:certificate}

The computation maintains an immutable set of canonical class keys for each
completed stage.  A key encodes the algebraic permutation of the canonical
cyclic-target representative, including algebraic fixed points.  Each stage
manifest records the format version, stage name, cyclic-domain size, number
of keys, and a SHA-256 checksum over deterministic key order.

Successful searches are stored separately as ordered triples of symbols.  A
witness checker reconstructs the cyclic-target pair, verifies that each
triple is applicable in the current cyclic order, applies it, checks that the
target cycle is preserved at every step, and finally checks that the algebraic
permutation is the identity.  Search failure is not encoded as a one-bit
claim: verification repeats the exact bounded search at the excluded depth.

The verification pass checks the following obligations.
\begin{enumerate}
  \item It recomputes every class-set count and checksum from the persistent
        key table.
  \item For $\En_{15}$, it checks absence of fixed points, repeats every
        depth-eight failure, and replays every depth-nine witness.
  \item For $\En_{16}$, it reconstructs all fixed-point expansions from
        $\En_{15}$, checks contraction of every stored fixed-point class, and
        repeats the depth-eight classification of the fixed-point-free part.
  \item For $\En_{17}$, it recomputes the complete canonical child set and the
        all-removals filter, repeats every depth-nine search, and replays every
        depth-ten witness.
  \item For order $18$, it recomputes the complete child set and removal
        filter and replays a depth-ten witness for every survivor.
\end{enumerate}
Generation is resumable at annotated-cycle-type boundaries.  Child
generation and classification use persistent processed-key sets, so an
interrupted pass may repeat work but cannot silently omit or duplicate a
class.  A completion manifest is written only after the corresponding class
set has been summarized.

The strongest final validation will use two reproducible passes from empty
storage and compare their per-stage key counts and checksums.  Small-order
tests independently enumerate ordinary permutations and confirm that the
quadratic insertion construction produces exactly the directly enumerated
child classes.  Further tests compare deletion multisets and child sets across
every representative in selected dihedral orbits.

\section{Computation status and remaining work}\label{sec:status}

Table~\ref{tab:status} distinguishes implemented machinery from completed
proof artifacts.  ``Implemented'' means that the construction and its
unit-level regression tests exist; it does not mean that the corresponding
astronomical enumeration has finished.

\begin{table}[ht]
\centering
\begin{tabular}{@{}lll@{}}
\toprule
Item & Software status & Proof-artifact status \\
\midrule
Class deletion and insertion & implemented and tested & small-order checks pass \\
Fixed-point-free direct generation & implemented and tested & large run pending \\
$\En_{15}$ classification & implemented & complete artifact pending \\
$\En_{16}$ classification & implemented & complete artifact pending \\
$\En_{17}$ propagation & implemented & depends on $\En_{16}$ \\
Order-$18$ certification & implemented & depends on $\En_{17}$ \\
Independent reproducibility pass & designed & pending \\
\bottomrule
\end{tabular}
\caption{Status of the working draft.}
\label{tab:status}
\end{table}

Before Theorem~\ref{thm:conditional} may be converted into an unconditional
diameter theorem, the computation must supply all stage manifests, the
complete witness stores, repeated exhaustive-search results, and agreement of
the independent passes.  The final version must also report per-route counts,
wall-clock and hardware information, peak storage, exact software revision,
and the location and checksum of the archived certificate.

\section{Conclusion}\label{sec:conclusion}

One-symbol deletion turns the order-$18$ diameter question into propagation
of a sequence of extremal class sets.  The key gain is structural: after the
complete $\En_{16}$ table has been computed, neither $S_{17}$ nor $S_{18}$ is
enumerated globally.  Only canonical children whose every deletion remains
extremal can advance.  Compatibility with extended-toric equivalence permits
this propagation directly at class level, using quadratically many raw
insertions per representative.

The resulting certificate has a sharp logical boundary.  Exhaustive failure
below the known diameter establishes extremal-set membership, while explicit
sorting sequences establish upper bounds.  Counts and checksums protect the
persisted class sets, and recomputation of child and removal filters protects
the inductive link between stages.  Completion of these obligations would
rule out distance eleven in $S_{18}$ and, with the lower bound, establish
$TD(18)=10$.

\bibliographystyle{plain}
\bibliography{references}

\end{document}
